\section{Blind Signature Construction}
\label{sec:blind-signature}

This section presents the vSIS-based blind signature that underlies ARC-SPRUCE. It combines the Ajtai commitment of
Section~\ref{subsec:commitments}, the trapdoor rational hash of Definition~\ref{def:rational-hash}, and the NIZK-ARK of
Definition~\ref{def:nizk-ark}. The signed message is always
\[
  \vecm=\vecmu\Vert \veck,
\]
where \(\vecmu\) denotes the attribute payload and \(\veck\) is a holder-generated secret PRF key used later for rate
limiting.

\subsubsection{Randomizing the rational-hash witness}
In the protocol, \(\Holder\) first commits to
\[
  \vecw=[\vecm\Vert \vecr_0^{\Holder}\Vert \vecr_1^{\Holder}\Vert \bar\vecr]^\top,
\]
where \((\vecr_0^{\Holder},\vecr_1^{\Holder})\) are the holder contributions to the rational-hash witness and
\(\bar\vecr\) is fresh commitment randomness. To satisfy the commitment bounds, we require
\(\|\vecw\|_\infty\le \BoundMsg\).

The issuance process then uses a two-party randomness generation step: the issuer samples
\((\vecr_0^{\Issuer},\vecr_1^{\Issuer})\), and the joint witness is obtained coefficient-wise by
\[
  \vecr_0=\centerfunc(\vecr_0^{\Holder}+\vecr_0^{\Issuer}),\qquad
  \vecr_1=\centerfunc(\vecr_1^{\Holder}+\vecr_1^{\Issuer}).
\]
Appendix~\ref{app:defs} records the centering map and its uniformity property.

\subsubsection{Blind Signature protocol}
\label{subsec:issuance-protocol}
We assume \(\nizkscheme=(\ZKProve,\ZKVerify)\) is a NIZK for the relations below,
\(\ComScheme=(\ComSetup,\Commit,\ComVerify)\) is the Ajtai commitment scheme, and
\[
  \TDRH=(\Setup,\TrapGen,\mathsf{Eval},\SampPre)
\]
is the trapdoor rational hash interface from Definition~\ref{def:bbs-rational-hash}. Public setup outputs the hash key
\(B\), while issuer key generation outputs the trapdoored signing matrix \(\mA\).

The pre-sign proof binds the commitment opening, the centering step, and the public target
\(\vect=h_{\vecm,(\vecr_0,\vecr_1)}(B)\). We write \(\relIssue(x,w)=1\) iff
\[
\begin{aligned}
  x&=\bigl(\ComPublicParam,\vecc,\vect,
  (\vecr_0^{\Issuer},\vecr_1^{\Issuer}),B\bigr),\\
  w&=(\vecm,\vecr_0^{\Holder},\vecr_1^{\Holder},\bar\vecr),\\
  \ComPublicParam&=(\mACom,\BoundMsg),\\
  \vecc&=\mACom\cdot[\vecm\Vert \vecr_0^{\Holder}\Vert \vecr_1^{\Holder}\Vert \bar\vecr]^\top,\\
  \|\vecm\|_\infty,\|\vecr_0^{\Holder}\|_\infty,\|\vecr_1^{\Holder}\|_\infty,\|\bar\vecr\|_\infty&\le \BoundMsg,\\
  \vecr_0&=\centerfunc(\vecr_0^{\Holder}+\vecr_0^{\Issuer}),\\
  \vecr_1&=\centerfunc(\vecr_1^{\Holder}+\vecr_1^{\Issuer}),\\
  \vect&=h_{\vecm,(\vecr_0,\vecr_1)}(B)\neq\bot.
\end{aligned}
\]

The issued blind signature is the short preimage witness together with the centered issuance randomness. We write
\(\relSign(x,w)=1\) iff
\[
\begin{aligned}
  x&=(\vecm,\mA,B,\BoundMsg,\BoundSig),\\
  w&=(\vecu,\vecr_0,\vecr_1),\\
  \|\vecr_0\|_\infty,\|\vecr_1\|_\infty&\le \BoundMsg,\qquad
  \|\vecu\|_\infty\le \BoundSig,\\
  h_{\vecm,(\vecr_0,\vecr_1)}(B)&\neq\bot,\qquad
  \mA\vecu=h_{\vecm,(\vecr_0,\vecr_1)}(B).
\end{aligned}
\]

We now define the blind-signature algorithms.
\begin{description}
  \item[$\Setup(\SecParam)$:] Fix the global ring parameters. Run
  \[
    (\mACom,\BoundMsg)\gets \ComScheme.\ComSetup(\SecParam)
  \]
  and
  \[
    (B,\mathcal{M},\mathcal{X},\delta,\BoundMsg,\BoundSig)\gets \TDRH.\Setup(\SecParam).
  \]
  Output
  \[
    \PublicParamsBS=(\mACom,B,\mathcal{M},\mathcal{X},\delta,\BoundMsg,\BoundSig).
  \]

  \item[$\IKeyGen()$:] Compute \((\mA,\trapdoor)\gets \TDRH.\TrapGen(\SecParam)\) and output
  \[
    \vk=\mA,\qquad \sk=\trapdoor.
  \]

  \item[$\Issue$:] The issuer \(\Issuer\) holds \(\sk\); the holder \(\Holder\) holds \(\vecmu,\vk\).
  \begin{enumerate}[leftmargin=1.25em]
    \item \(\Holder\) samples \(\veck,\vecr_0^{\Holder},\vecr_1^{\Holder},\bar\vecr\) coefficient-wise in
    \([-\BoundMsg,\BoundMsg]\) and sets \(\vecm=\vecmu\Vert \veck\).
    \item \(\Holder\) computes
    \[
      \vecc\gets \mACom\cdot[\vecm\Vert \vecr_0^{\Holder}\Vert \vecr_1^{\Holder}\Vert \bar\vecr]^\top
    \]
    and sends \(\vecc\) to \(\Issuer\).
    \item \(\Issuer\) samples \(\vecr_0^{\Issuer},\vecr_1^{\Issuer}\) coefficient-wise in
    \([-\BoundMsg,\BoundMsg]\) and sends them to \(\Holder\).
    \item \(\Holder\) computes
    \[
      \vecr_0\gets \centerfunc(\vecr_0^{\Holder}+\vecr_0^{\Issuer}),\qquad
      \vecr_1\gets \centerfunc(\vecr_1^{\Holder}+\vecr_1^{\Issuer}),
    \]
    then evaluates \(\vect\gets h_{\vecm,(\vecr_0,\vecr_1)}(B)\). If \(\vect=\bot\), it aborts and restarts with fresh
    local randomness.
    \item \(\Holder\) computes \(\nizkIssue\gets \nizkscheme.\ZKProve(w,x)\) for the relation \(\relIssue\), with
    witness \(w=(\vecm,\vecr_0^{\Holder},\vecr_1^{\Holder},\bar\vecr)\) and statement
    \(x=(\ComPublicParam,\vecc,\vect,\vecr_0^{\Issuer},\vecr_1^{\Issuer},B)\), and sends \((\vect,\nizkIssue)\) to
    \(\Issuer\).
    \item \(\Issuer\) verifies \(\nizkIssue\) for \(\relIssue\). If verification fails, it aborts. Otherwise it computes
    \[
      \vecu\gets \TDRH.\SampPre(\trapdoor,\vect)
    \]
    and sends \(\vecu\) to \(\Holder\).
    \item \(\Holder\) checks \(\|\vecu\|_\infty\le \BoundSig\) and \(\mA\vecu=\vect\). If so, it outputs
    \[
      \sig=(\vecu,\vecr_0,\vecr_1)
    \]
    on message \(\vecm=\vecmu\Vert \veck\); otherwise it aborts.
  \end{enumerate}

  \item[$\Verify(\vk,\vecm,\sig)$:] Let \(\vk=\mA\). Parse \(\sig=(\vecu,\vecr_0,\vecr_1)\). Accept iff
  \[
    \bigl((\vecm,\mA,B,\BoundMsg,\BoundSig),(\vecu,\vecr_0,\vecr_1)\bigr)\in \relSign.
  \]
\end{description}

For ARC, the holder does not reveal \(\sig\) directly. Section~\ref{sec:arc-construction} instead proves the relation
\(\relSign\) in zero knowledge while keeping \(\veck,\vecr_0,\vecr_1\) hidden and binding \(\veck\) to the PRF tag.

\subsubsection{Security}
\begin{proposition}[Correctness]
Assume the pre-sign NIZK is complete and the trapdoor sampler is correct. Then an honest execution of \(\Issue\)
outputs a signature \(\sig\) that is accepted by \(\Verify\), except with negligible failure probability coming from the
trapdoor sampler or an inadmissible hash evaluation.
\end{proposition}

\begin{proof}
Completeness of the pre-sign NIZK implies that the issuer accepts the honest relation \(\relIssue\). By correctness of
\(\SampPre\), the returned \(\vecu\) satisfies \(\mA\vecu=\vect\) and the public shortness bound except with negligible
probability. The output \(\sig=(\vecu,\vecr_0,\vecr_1)\) therefore satisfies \(\relSign\).
\end{proof}

\begin{remark}[Security roadmap]
The non-blind hash-and-preimage signature layer underlying this protocol is justified by the DKLW25 reduction discussed
in Remark~\ref{rem:vsis-unforgeability}. For the blind protocol itself, the intended route is to combine straightline
extractability of the pre-sign proof with an interactive \(\mathsf{GenISIS}_f\)-style assumption for
\[
  f(B,\vecm,(\vecr_0,\vecr_1))=h_{\vecm,(\vecr_0,\vecr_1)}(B),
\]
while using statistical hiding of the Ajtai commitment and zero knowledge of the pre-sign proof in the blindness
argument. This explains why the construction targets one-more unforgeability and blindness, but we do not claim a
completed reduction for either property here.
\end{remark}
